% Education
% =================================================
\DeclareTranslation{portuguese}{section-education}{FORMAÇÃO}

% ENSTA
\DeclareTranslation{portuguese}{ensta-title}{Mestrado em Engenharia de Robótica}
\DeclareTranslation{portuguese}{ensta-site}{École Nationale Supérieure de Techniques Avancées (ENSTA), Institut Polytechnique de Paris (IPP)}
\DeclareTranslation{portuguese}{ensta-location}{Palaiseau, Île-de-France, França}
\DeclareTranslation{portuguese}{ensta-duration}{jul 2022 -- out 2025}
\DeclareTranslation{portuguese}{ensta-description}{%
CR: 3,71/4,00. Programa de duplo diploma de três anos com a UNICAMP.
}

% UNICAMP
\DeclareTranslation{portuguese}{unicamp-title}{Bacharelado em Engenharia de Controle e Automação}
\DeclareTranslation{portuguese}{unicamp-site}{Faculdade de Engenharia Mecânica (FEM), Universidade Estadual de Campinas (UNICAMP)}
\DeclareTranslation{portuguese}{unicamp-location}{Campinas, São Paulo, Brasil}
\DeclareTranslation{portuguese}{unicamp-duration}{fev 2019 -- ago 2026}
\DeclareTranslation{portuguese}{unicamp-description}{%
CR: 0,84/1,00. Classificado em 10º entre 59 estudantes e recebeu a Láurea Acadêmica pelo desempenho acadêmico e conclusão no prazo.
}

% Experiences
% =================================================
\DeclareTranslation{portuguese}{section-experiences}{EXPERIÊNCIA}

% Ansys
\DeclareTranslation{portuguese}{ansys-title}{Estagiário de Engenharia de Garantia da Qualidade}
\DeclareTranslation{portuguese}{ansys-site}{ANSYS, Synopsys} % codespell:ignore
\DeclareTranslation{portuguese}{ansys-location}{La Farlède, Provence-Alpes Côte d'Azur, França}
\DeclareTranslation{portuguese}{ansys-duration}{mai 2025 -- out 2025}
\DeclareTranslation{portuguese}{ansys-description}{%
\begin{itemize}[nosep, leftmargin=*]
\item \textbf{Automação de Testes:} Automatizei 43 testes manuais de regressão em workflows do Azure Pipelines usando IronPython para a suíte de simulação óptica Speos, cobrindo cenários de sensores LiDAR e modelagem de lentes. O workflow comparava os resultados e gerava relatórios de aprovação/reprovação, eliminando mais de 1.200 minutos de testes manuais por versão.
\item \textbf{Visão Computacional:} Desenvolvi do zero uma ferramenta configurável de comparação de imagens em Python usando OpenCV e scikit-image para testes de regressão, combinando SSIM, detecção de desfoque baseada em Laplaciano, análise de histogramas de cores.
\end{itemize}
}
\DeclareTranslation{portuguese}{ansys-skills}{Azure Pipelines, CI/CD, Git, GitHub Actions, gRPC, IronPython, OpenCV, PyTest, Python, scikit-image}

% SYSNAV
\DeclareTranslation{portuguese}{sysnav-title}{Estagiário de Firmware Embarcado e Hardware}
\DeclareTranslation{portuguese}{sysnav-site}{SYSNAV}
\DeclareTranslation{portuguese}{sysnav-location}{Vernon, Normandia, França}
\DeclareTranslation{portuguese}{sysnav-duration}{mai 2024 -- set 2024}
\DeclareTranslation{portuguese}{sysnav-description}{%
\begin{itemize}[nosep, leftmargin=*]
\item \textbf{Desenvolvimento de Firmware de Baixo Nível:} Desenvolvi do zero, em C, um driver de firmware I2C para um circuito integrado carregador de bateria da Texas Instruments em uma base de código STM32H7. Projetei a máquina de estados finitos do driver e validei a proteção contra sobrecorrente, a proteção térmica e as transições de estado de carga com base no datasheet do dispositivo.
\item \textbf{Testes Hardware-in-the-Loop:} Desenvolvi um framework de testes hardware-in-the-loop em Python usando uma placa STM32H7 Nucleo e 16 relés controlados por I2C para simular falhas de conexão entre a estação, a base e o relógio do dispositivo SYDE. Projetei a biblioteca de orquestração em Python e automatizei o controle dos relés e a recuperação dos logs do dispositivo. Reduzi o tempo de validação do firmware de aproximadamente 2 dias para 2 horas, reproduzindo falhas de campo, como falhas na detecção do relógio.
\end{itemize}
}
\DeclareTranslation{portuguese}{sysnav-skills}{C, Git, GitLab, Hardware-in-the-Loop, I2C, Analisador Lógico, Osciloscópio, Python, STM32H7}

% Amazon
\DeclareTranslation{portuguese}{amazon-title}{Estagiário de Gestão de Fornecedores de Varejo}
\DeclareTranslation{portuguese}{amazon-site}{Amazon}
\DeclareTranslation{portuguese}{amazon-location}{Clichy, Île-de-France, França}
\DeclareTranslation{portuguese}{amazon-duration}{set 2023 -- mar 2024}
\DeclareTranslation{portuguese}{amazon-description}{%
\begin{itemize}[nosep, leftmargin=*]
\item \textbf{Automação de Processamento de Dados:} Desenvolvi uma ferramenta de correspondência de palavras-chave em VBA para classificar dados de licenciamento de marcas em mais de 400.000 registros, reduzindo o tempo de processamento de 6 horas para 20 minutos.
\end{itemize}
}
\DeclareTranslation{portuguese}{amazon-skills}{Business Intelligence, Análise de Dados, Microsoft Excel, Otimização de Processos, SQL, VBA}

% ISIR
\DeclareTranslation{portuguese}{isir-title}{Estagiário de Pesquisa em Neurociência}
\DeclareTranslation{portuguese}{isir-site}{ISIR, Sorbonne Université}
\DeclareTranslation{portuguese}{isir-location}{Paris, Île-de-France, França}
\DeclareTranslation{portuguese}{isir-duration}{mai 2023 -- ago 2023}
\DeclareTranslation{portuguese}{isir-description}{%
\begin{itemize}[nosep, leftmargin=*]
\item \textbf{Simulação Sensório-Motora:} Implementei um ambiente de simulação em Python/Webots baseado em modelos matemáticos de cognição espacial do hipocampo, incluindo análise de imagens com OpenCV para mapeamento da navegação.
\end{itemize}
}
\DeclareTranslation{portuguese}{isir-skills}{Visualização de Dados, Git, LaTeX, Matplotlib, OpenCV, Python, Simulação Robótica, SciPy, Modelagem Estatística, Webots}

% Phoenix
\DeclareTranslation{portuguese}{phoenix-title}{Líder de Eletrônica}
\DeclareTranslation{portuguese}{phoenix-site}{Equipe Phoenix de Robótica da UNICAMP}
\DeclareTranslation{portuguese}{phoenix-location}{Campinas, São Paulo, Brasil}
\DeclareTranslation{portuguese}{phoenix-duration}{jul 2020 -- jul 2021}
\DeclareTranslation{portuguese}{phoenix-description}{%
\begin{itemize}[nosep, leftmargin=*]
\item \textbf{Liderança:} Liderei uma divisão de eletrônica com 12 membros, desenvolvendo treinamentos de Altium e eletrônica para novos integrantes e coordenando o desenvolvimento, fabricação e diagnóstico de placas.
\end{itemize}
}
\DeclareTranslation{portuguese}{phoenix-skills}{Altium 365, Altium Designer, Diagnóstico de Hardware, Projeto de PCB, Fabricação e Montagem.}

% Languages
% =================================================
\DeclareTranslation{portuguese}{section-languages}{IDIOMAS}

% english
\DeclareTranslation{portuguese}{english-name}{Inglês}
\DeclareTranslation{portuguese}{english-cefr}{C1}
\DeclareTranslation{portuguese}{english-hbar}{83}%
\DeclareTranslation{portuguese}{english-level}{Fluente}

% french
\DeclareTranslation{portuguese}{french-name}{Francês}
\DeclareTranslation{portuguese}{french-cefr}{C1}
\DeclareTranslation{portuguese}{french-hbar}{83}%
\DeclareTranslation{portuguese}{french-level}{Fluente}

% Portuguese
\DeclareTranslation{portuguese}{portuguese-name}{Português}
\DeclareTranslation{portuguese}{portuguese-cefr}{C2}
\DeclareTranslation{portuguese}{portuguese-hbar}{100}%
\DeclareTranslation{portuguese}{portuguese-level}{Nativo}

% Spanish
\DeclareTranslation{portuguese}{spanish-name}{Espanhol}
\DeclareTranslation{portuguese}{spanish-cefr}{B1}
\DeclareTranslation{portuguese}{spanish-hbar}{50}%
\DeclareTranslation{portuguese}{spanish-level}{Intermediário}

% Presentation
% =================================================
\DeclareTranslation{portuguese}{presentation}{%
Engenheiro de Software Embarcado com experiência prática no desenvolvimento de firmware em C para microcontroladores STM32 e ESP32-S3, incluindo drivers de periféricos, protocolos de comunicação e máquinas de estados finitos. Experiência em validação hardware-in-the-loop, diagnóstico de hardware, CI/CD e desenvolvimento de sistemas robóticos. Mestre em Engenharia de Robótica pela ENSTA Paris.
}

% Projects
% =================================================
\DeclareTranslation{portuguese}{section-project}{PROJETOS}

% dolphin
\DeclareTranslation{portuguese}{dolphin-title}{Desenvolvedor de Software Embarcado e Hardware}
\DeclareTranslation{portuguese}{dolphin-site}{Equipe Paralela -- Dolphin, um controlador de robô de sumô de código aberto}
\DeclareTranslation{portuguese}{dolphin-location}{Campinas, São Paulo, Brasil}
\DeclareTranslation{portuguese}{dolphin-duration}{fev 2026 -- presente}
\DeclareTranslation{portuguese}{dolphin-description}{%
\begin{itemize}[nosep, leftmargin=*]
\item \textbf{Firmware:} Desenvolvendo firmware em C para um controlador de robô ESP32-S3 usando ESP-IDF, PlatformIO, CMake e uma camada mínima de FreeRTOS, implementando o controle por máquina de estados finitos orientada por tabelas para operação autônoma e controlada por rádio. Validei o controlador em um robô de competição funcional, terminando em 6º lugar entre 11 robôs na competição mais recente.
\item \textbf{Hardware:} Projetei a placa principal do controlador, o sensor de linha QRE1113 e as placas do encoder AS5147U em KiCAD, integrando controle PWM de motores, aquisição de sensores por ADC, multiplexação controlada por GPIO, periféricos SPI e receptores RC e sensores infravermelhos acionados por interrupções.
\item \textbf{Infraestrutura:} Desenvolvi CI com GitHub Actions para compilação do firmware e validação ERC/DRC do KiCAD.
\end{itemize}
}
\DeclareTranslation{portuguese}{dolphin-skills}{C, ESP32-S3, Git, GitHub Actions, KiCAD, PlatformIO.}
